%% ---------------------------------------------------------------
%%  Z_p-valued residues for regular singular stratified bundles
%%  Working draft, version 2.
%%  Changes from v1:
%%    * Section 8 rewritten; the relative-structure GAP is closed.
%%    * New Section 9: p-adic middle convolution.
%% ---------------------------------------------------------------
\documentclass[11pt,reqno]{amsart}

\usepackage{amsmath,amssymb,amsthm,mathrsfs}
\usepackage[T1]{fontenc}
\usepackage{microtype}
\usepackage{stmaryrd}
\usepackage[colorlinks,linkcolor=blue!60!black,citecolor=green!40!black]{hyperref}
\usepackage[capitalise]{cleveref}
\usepackage{xcolor}

%% ---- editorial macros (remove before submission) ---------------
\newcommand{\TODO}[1]{{\color{red}\textsf{[TODO: #1]}}}
\newcommand{\GAP}[1]{{\color{orange}\textsf{[GAP: #1]}}}
\newcommand{\CHECK}[1]{{\color{blue!70!black}\textsf{[check: #1]}}}

%% ---- notation --------------------------------------------------
\newcommand{\Zp}{\mathbb{Z}_p}
\newcommand{\Fp}{\mathbb{F}_p}
\newcommand{\ZZ}{\mathbb{Z}}
\newcommand{\QQ}{\mathbb{Q}}
\newcommand{\NN}{\mathbb{N}}
\newcommand{\GG}{\mathbb{G}}
\newcommand{\PP}{\mathbb{P}}
\newcommand{\AAA}{\mathbb{A}}
\newcommand{\cO}{\mathcal{O}}
\newcommand{\cD}{\mathcal{D}}
\newcommand{\cE}{\mathcal{E}}
\newcommand{\cF}{\mathcal{F}}
\newcommand{\cG}{\mathcal{G}}
\newcommand{\cH}{\mathcal{H}}
\newcommand{\cL}{\mathcal{L}}
\newcommand{\cM}{\mathcal{M}}
\newcommand{\cA}{\mathcal{A}}
\newcommand{\cR}{\mathcal{R}}
\newcommand{\cT}{\mathcal{T}}
\newcommand{\fg}{\mathfrak{g}}
\newcommand{\ft}{\mathfrak{t}}
\newcommand{\Strat}{\mathrm{Strat}}
\newcommand{\rs}{\mathrm{rs}}
\newcommand{\tame}{\mathrm{tame}}
\newcommand{\strat}{\mathrm{strat}}
\newcommand{\Pic}{\mathrm{Pic}}
\newcommand{\Vect}{\mathrm{Vect}}
\newcommand{\Rep}{\mathrm{Rep}}
\newcommand{\Hom}{\mathrm{Hom}}
\newcommand{\End}{\mathrm{End}}
\newcommand{\Ext}{\mathrm{Ext}}
\newcommand{\Spec}{\mathrm{Spec}}
\newcommand{\Frac}{\mathrm{Frac}}
\newcommand{\res}{\mathrm{res}}
\newcommand{\inv}{\mathrm{inv}}
\newcommand{\gr}{\mathrm{gr}}
\newcommand{\Iw}{\mathrm{Iw}}
\newcommand{\dig}{\mathrm{dig}}
\newcommand{\rk}{\mathrm{rk}}
\newcommand{\id}{\mathrm{id}}
\newcommand{\im}{\mathrm{im}}
\newcommand{\tr}{\mathrm{tr}}
\newcommand{\MC}{\mathrm{MC}}
\newcommand{\rig}{\mathrm{rig}}
\newcommand{\Xst}{X_*(T)}
\newcommand{\isoto}{\xrightarrow{\ \sim\ }}
\newcommand{\dR}{\mathrm{dR}}

%% ---- theorem environments --------------------------------------
\theoremstyle{plain}
\newtheorem{theorem}{Theorem}[section]
\newtheorem{proposition}[theorem]{Proposition}
\newtheorem{lemma}[theorem]{Lemma}
\newtheorem{corollary}[theorem]{Corollary}
\newtheorem{conjecture}[theorem]{Conjecture}
\newtheorem{problem}[theorem]{Problem}

\theoremstyle{definition}
\newtheorem{definition}[theorem]{Definition}
\newtheorem{example}[theorem]{Example}
\newtheorem{construction}[theorem]{Construction}

\theoremstyle{remark}
\newtheorem{remark}[theorem]{Remark}
\newtheorem{notation}[theorem]{Notation}

\numberwithin{equation}{section}

%% ---------------------------------------------------------------
\begin{document}

\title[$\Zp$-valued residues]{$\Zp$-valued residues for regular singular
stratified bundles, boundary restriction, and $p$-adic middle convolution}

\author{}
\address{}
\email{}

\date{\today\\ \textsc{Working draft v2 --- not for circulation}}

\begin{abstract}
Let $k$ be algebraically closed of characteristic $p>0$, $(\bar X,D)$ a
smooth $k$-variety with normal crossings divisor, $U=\bar X\setminus D$.
We attach to every regular singular stratified bundle on $U$, together
with a logarithmic lattice, a \emph{$\Zp$-valued residue} along each
boundary component: the level-$m$ logarithmic residues lie in $\Fp$ and
are exactly the $p$-adic digits of a single element of $\Xst\otimes\Zp$,
well defined modulo $\Xst$.

Our local result is that $\Strat^{\rs}(k((t)))$ is semisimple and
equivalent to $\Rep\bigl(\Spec k[\Zp/\ZZ]\bigr)$: local monodromy is
always semisimple and the $\Zp$-residue is a complete invariant.  We then
prove a boundary restriction theorem: $\bar E_0|_{D_i}$ carries a
canonical decomposition indexed by $\Zp/\ZZ$, each piece being a
regular singular stratified bundle on $D_i$ \emph{twisted by
$N_i^{\otimes c}$}, $N_i$ the normal bundle.  The obstruction to
untwisting lies in $\widehat{\Pic(D_i)}_p/\Pic(D_i)$, and vanishing of the
residue is exactly the criterion for extending across $D_i$.

Finally we construct a $p$-adic middle convolution $\MC_c$ for
$c\in\Zp$ on $\PP^1$.  Because local monodromy is semisimple, Katz's
Jordan-block bookkeeping disappears entirely and the effect of $\MC_c$ on
the $\Zp$-valued Riemann scheme is purely additive.  We formulate the
resulting $p$-adic Katz algorithm and define $p$-adic hypergeometric
stratified bundles with parameters in $\Zp$.
\end{abstract}

\maketitle

\setcounter{tocdepth}{1}
\tableofcontents

%% ===============================================================
\section{Introduction}
%% ===============================================================

\subsection{}
In characteristic zero a regular singular connection has a residue along
the boundary whose eigenvalues modulo $\ZZ$ record local monodromy at
infinity.  In characteristic $p>0$ the analogue of a flat connection is a
stratified bundle (Katz \cite{Katz}), i.e.\ an $F$-divided sheaf
$\cE=\bigl((E_n)_{n\ge0},\sigma_n\colon F^*E_{n+1}\isoto E_n\bigr)$, and
regular singularity was defined by Kindler \cite{Kindler}.  A stratified
bundle carries a residue at \emph{each level}, and all of them lie in
$\Fp$.  The purpose of this paper is to bundle them into one $p$-adic
number and to develop the resulting theory.

\subsection{The mechanism}
In a local coordinate $t$ cutting out $D$ set
\[
   D_a:=t^{a}\partial_t^{[a]}\in\cD_{\bar X}(\log D),\qquad a\ge0,
   \qquad \delta^{(m)}:=D_{p^m}.
\]
These act on $t^n$ by $\binom na$, hence commute and satisfy $D_a^p=D_a$;
consequently \emph{all} logarithmic residues are semisimple with
eigenvalues in $\Fp$ (\Cref{cor:rigidity}).  Since
$\binom{\cdot}{p^m}\colon\Zp\to\Fp$ is continuous and returns the $m$-th
$p$-adic digit (Lucas), a compatible family of level residues is precisely
the digit expansion of a single $r\in\Zp$.  What makes this more than
repackaging is that the \emph{ring} structure of $\Zp$ matches the
geometry: eigenvalues add under $\otimes$, multiply by $m$ under tame
covers of degree $m$, and satisfy $\sum_x r_x=-\deg$.

\subsection{Main results}
\begin{itemize}
\item \Cref{thm:localstructure}: $\Strat^{\rs}(K)\simeq\Rep(\cT_p)$,
      $\cT_p=\Spec k[\Zp/\ZZ]$; local monodromy is semisimple and $r$ is
      complete.  In particular there is no local unipotent part
      (\Cref{cor:nounipotent}).
\item \Cref{thm:comparison}: $r=-\lim\nu_n^{\Iw}$ for the flag-wise
      (Iwasawa) Hecke invariants; the naive Cartan invariants fail
      (\Cref{ex:cartanfails}).
\item \Cref{thm:boundary}: canonical $\Zp/\ZZ$-graded boundary
      restriction, with twist $N_i^{\otimes c}$; obstruction in
      $\widehat{\Pic(D_i)}_p/\Pic(D_i)$; extension criterion
      \Cref{cor:extension}.
\item \Cref{sec:MC}: construction of $\MC_c$ for $c\in\Zp$, the exponent
      rules (\Cref{prop:MCrules}), invariance of the rigidity index, and
      the $p$-adic Katz algorithm (\Cref{conj:katzalg}).
\end{itemize}

\subsection{Changes from the first draft}
Two corrections are worth recording, since both were instructive.
\begin{enumerate}
\item The level residues are the digits of $r$, but the digits do
      \emph{not} add: $\dig_m(c)+\dig_m(c')\ne\dig_m(c+c')$.  All
      additivity statements must be made for the eigenvalues in $\Zp$, and
      the proof runs through Vandermonde's identity for the full family
      $\{D_a\}$, not just $\{\delta^{(m)}\}$ (\Cref{thm:functoriality}).
\item In \Cref{thm:boundary} the twisting exponent at level $n$ is the
      \emph{Hecke exponent} $a_n=pc_{n+1}-c_n\in\ZZ$, not the digit
      $\dig_n(c)$.  The two agree only after normalising the lattices;
      the telescoped total twist $N_i^{\otimes c}$ is the same either way.
\end{enumerate}

%% ===============================================================
\section{Conventions}\label{sec:conventions}
%% ===============================================================

$k$ is algebraically closed of characteristic $p>0$, $F$ absolute
Frobenius, $\bar X$ smooth over $k$, $D=\sum_{i\in I}D_i$ a strict normal
crossings divisor with smooth components, $U=\bar X\setminus D$.  Locally
$R=k[[t]]$, $K=k((t))$.

\begin{definition}
A \emph{stratified bundle} on $U$ is
$\cE=\bigl((E_n)_{n\ge0},\sigma_n\bigr)$ with $E_n\in\Vect(U)$ and
$\sigma_n\colon F^*E_{n+1}\isoto E_n$.  Write $E:=E_0$; the isomorphism
$E\cong F^{n*}E_n$ makes $E$ a module over the subring
$\cD_U^{(n)}\subset\cD_U$ generated by operators of order $<p^n$, and
$\cE$ is the same as an $\cO$-locally free $\cD_U$-module.  $\Strat(U)$ is
tannakian over $k$.
\end{definition}

\begin{definition}\label{def:lattice}
A \emph{lattice system} is a family $\bar E_n\in\Vect(\bar X)$ with
$\bar E_n|_U=E_n$; then
$\bar\sigma_n\colon F^*\bar E_{n+1}\to\bar E_n$ is a modification along
$D$.  It is \emph{logarithmic} if each $\bar E_n$ is
$\cD_{\bar X}(\log D)$-stable.  $\cE$ is \emph{regular singular} if a
logarithmic lattice system exists (equivalently \cite{Kindler}).
\end{definition}

\begin{notation}\label{not:Da}
$D_a:=t^a\partial_t^{[a]}$, $\delta^{(m)}:=D_{p^m}$,
$\theta_a:=\nabla(D_a)\in\End_k(\bar E_0)$,
$R_a:=\res_D\theta_a\in\End_{\cO_D}(\bar E_0|_D)$, $R^{(m)}:=R_{p^m}$.
For $c\in\Zp$, $\dig_m(c)\in\{0,\dots,p-1\}$ is the $m$-th digit.
\end{notation}

\begin{definition}\label{def:logpar}
For $G$ split reductive over $\Fp$, $P\subset G$ parabolic, $T\subset P$
maximal torus, $W$ Weyl group, a \emph{log parabolic geometry of type
$(G,P)$} on $(\bar X,D)$ is a $P$-bundle $\cG$ with $P$-equivariant
$\omega\colon T\cG(-\log\tilde D)\isoto\fg\otimes\cO$.  An
\emph{infinite-level} such geometry is a family $(\cG_n,\omega_n)$ whose
adjoint tractor bundles $\cA_n=\cG_n\times^P\fg$ form a logarithmic
lattice system.  All residue statements below concern $\cA_\bullet$ alone.
\end{definition}

%% ===============================================================
\section{Rigidity of logarithmic residues}\label{sec:rigidity}
%% ===============================================================

\begin{proposition}\label{prop:idempotent}
In $\cD_{\bar X}(\log D)$, for all $a,b\ge0$:
\[
   D_a\,t^n=\binom na\,t^n,\qquad D_a^{\,p}=D_a,\qquad [D_a,D_b]=0,
   \qquad D_a=\prod_{m\ge0}\binom{D_{p^m}}{a_m},
\]
the last identity for $a=\sum_ma_mp^m$ in base $p$.
\end{proposition}
\begin{proof}
$\cD_{\bar X}(\log D)\subset\End_k(\cO)$ is faithful and
$\partial^{[a]}t^n=\binom nat^{n-a}$ gives the first formula; the $D_a$ are
therefore simultaneously diagonal in the basis $\{t^n\}$, whence
commutativity, and $\binom na\in\Fp$ satisfies $c^p=c$, whence
$D_a^p=D_a$.  The last identity is Lucas' congruence
$\binom na\equiv\prod_m\binom{n_m}{a_m}$ evaluated on the same basis.
\end{proof}

\begin{corollary}\label{cor:rigidity}
The $\theta_a$ are commuting $k$-linear operators with
$\theta_a^p=\theta_a$ and $\theta_a=\prod_m\binom{\theta_{p^m}}{a_m}$.
The residues $R_a$ are simultaneously diagonalisable with eigenvalues in
$\Fp$; their joint characters are of the form
$a\mapsto\binom{c}{a}$ for a unique $c\in\Zp$.
\end{corollary}
\begin{proof}
Everything but the last clause is immediate.  A joint character is
determined by its values on $a=p^m$ by the Lucas identity; choosing
$c\in\Zp$ with $\dig_m(c)$ equal to the $\delta^{(m)}$-eigenvalue and
invoking \Cref{lem:digitmap} gives $\binom{c}{a}$ for all $a$.
\end{proof}

\begin{lemma}[Leibniz]\label{lem:leibniz}
$\theta_a(fx)=\sum_{b+c=a}D_b(f)\theta_c(x)$; in particular
$\theta_a(tx)=t\theta_a(x)+t\theta_{a-1}(x)$, so each $\theta_a$ preserves
$t\bar E_0$ and $R_a$ is defined and $\cO_D$-linear.
\end{lemma}

\begin{lemma}\label{lem:digitmap}
$x\mapsto\binom{x}{p^m}\bmod p$ extends uniquely to a continuous map
$\Zp\to\Fp$, equal to $\dig_m$.
\end{lemma}
\begin{proof}
By Lucas, $\binom{n}{p^m}\bmod p$ depends only on $n\bmod p^{m+1}$ and
equals $\dig_m(n)$ for $n\ge0$; extend by density.
\end{proof}

\begin{lemma}\label{lem:constant}
If $D_i$ is connected the characteristic polynomial of $R_a$ is constant
on $D_i$ and splits over $\Fp$.
\end{lemma}
\begin{proof}
Eigenvalues lie in the discrete set $\Fp$, so multiplicities are locally
constant.
\end{proof}

%% ===============================================================
\section{The $\Zp$-valued residue}\label{sec:residue}
%% ===============================================================

\begin{definition}[$r$-spectrum]\label{def:rspec}
Let $\bar E_0$ be a logarithmic lattice.  By \Cref{cor:rigidity},
\[
   \bar E_0|_{D_i}=\bigoplus_{c\in\Zp}\bigl(\bar E_0|_{D_i}\bigr)^c,
   \qquad
   \bigl(\bar E_0|_{D_i}\bigr)^c=\bigl\{x:\ R_ax=\tbinom ca x\ \forall a\bigr\},
\]
a finite direct sum of $\cO_{D_i}$-subbundles.  The \emph{$r$-spectrum}
is the multiset of occurring $c$ with multiplicities; for a log parabolic
geometry it is an element
$r_i\in\bigl(\Xst\otimes\Zp\bigr)/W$, the \emph{$\Zp$-valued residue}.
\end{definition}

\begin{theorem}\label{thm:welldef}
\begin{enumerate}
\item The $r$-spectrum depends only on $\bar E_0$, not on
      $\bar E_n$ for $n\ge1$.
\item A Hecke modification of $\bar E_0$ of type $\nu\in\Xst$ changes $r$
      to $r+\nu$.
\item Hence $\bar r_i\in\bigl(\Xst\otimes\Zp/\Xst\bigr)/W$ depends only on
      $\cE$.
\end{enumerate}
\end{theorem}
\begin{proof}
(1) is immediate from the definition.  (2) is \Cref{prop:jump} below.
\end{proof}

\begin{theorem}\label{thm:functoriality}
\begin{enumerate}
\item\label{it:additive} $\mathrm{Spec}_r(\cE\otimes\cE')
      =\{c+c'\}$, sums in $\Zp$.
\item $\mathrm{Spec}_r(\cE^\vee)=-\mathrm{Spec}_r(\cE)$.
\item For $\phi\colon G\to G'$, $r'=\mathrm{Lie}(\phi)(r)$.
\item\label{it:pullback} If $f^*D_i=\sum_je_{ij}E_j$ then
      $r_{E_j}(f^*\cE)=\sum_ie_{ij}r_{D_i}(\cE)$.
\end{enumerate}
\end{theorem}
\begin{proof}
\eqref{it:additive}: on simultaneous eigenvectors, \Cref{lem:leibniz}
applied to $\cE\otimes\cE'$ gives
$\theta_a(x\otimes x')=\sum_{b+c=a}\theta_b x\otimes\theta_c x'$, so the
eigenvalue is $\sum_{b+c=a}\binom{c}{b}\binom{c'}{c}=\binom{c+c'}{a}$ by
Vandermonde.  The rest follows.
\end{proof}

\begin{remark}
It is essential that eigenvalues add \emph{in $\Zp$}: digits do not add,
because of carries.  This is the first place where the ring structure of
$\Zp$, rather than the family of digits, does real work.
\end{remark}

\begin{corollary}\label{cor:tame}
For a Kummer cover of ramification index $m$, $p\nmid m$: $r\mapsto mr$.
\end{corollary}

%% ===============================================================
\section{The local structure theorem}\label{sec:localstructure}
%% ===============================================================

Here $R=k[[t]]$, $K=k((t))$.  For $c\in\Zp$ let $\cL_c$ be the rank one
stratified bundle with $\sigma_n(1\otimes e_{n+1})=t^{-\dig_n(c)}e_n$; then
$\mathrm{Spec}_r(\cL_c)=\{c\}$.

\begin{lemma}\label{lem:mu0}
If every $\sigma_n$ is an isomorphism of lattices then $\cE$ is trivial.
\end{lemma}
\begin{proof}
Let $g_n\in GL_N(R)$ represent $\sigma_n$; trivialising means solving
$v_n=v_{n+1}^{(p)}g_n$, where $(\cdot)^{(p)}$ applies $F$ entrywise.
Reducing mod $t$ gives an $F$-divided sheaf on $\Spec k$, trivial by
Lang's theorem; so we may assume $g_n\in1+tM_N(R)$.  Then
$g_{n+i}^{(p^i)}\in1+t^{p^i}M_N(R)$ and
$v_n:=\lim_Ng_{n+N}^{(p^N)}\cdots g_n$ converges $t$-adically in
$1+tM_N(R)$ and solves the equation.
\end{proof}

\begin{lemma}\label{lem:res0}
If $R^{(0)}=0$ then $\bar E_1$ can be replaced so that
$\sigma_0(F^*\bar E_1)=\bar E_0$.
\end{lemma}
\begin{proof}
$\theta_1$ is semisimple with eigenvalues in $\Fp$; write
$\bar E_0=\bigoplus_{c\in\Fp}\bar E_0^{(c)}$.  By \Cref{lem:leibniz},
\begin{equation}\label{eq:grading}
   t\cdot\bar E_0^{(c)}\subseteq\bar E_0^{(c+1)},
\end{equation}
i.e.\ $\bar E_0$ is $\Fp$-graded with $\deg t=1$.  The hypothesis says
$\theta_1(\bar E_0)\subseteq t\bar E_0$; for $c\ne0$ and
$x\in\bar E_0^{(c)}$ this forces $x\in t\bar E_0$.  So
$\bar E_0=M+t\bar E_0$ with $M:=\bar E_0^{(0)}$, hence $\bar E_0=R\cdot M$
by completeness.  By \eqref{eq:grading}, $t^pM\subseteq M$, the sum
$\sum_{j<p}t^jM$ is direct and equals $\bar E_0$, so $M$ is free of rank
$N$ over $k[[t^p]]$.  Finally $\theta_1x=0$ gives $\nabla(\partial)x=0$ and
hence $\nabla(\partial^{[j]})x=\nabla(\partial)^jx/j!=0$ for $j<p$, so $M$
is level-$1$ horizontal; Cartier descent produces $\bar E_1$.
\end{proof}

\begin{lemma}\label{lem:shift}
If $\sigma_0$ is a lattice isomorphism then $\bar E_1$ is
$\cD(\log)$-stable and $r(\cE)=p\cdot r(\cE_1)$, where
$\cE_1=(E_{1+i},\sigma_{1+i})$.
\end{lemma}

\begin{corollary}\label{cor:allzero}
$r=0$ implies $\cE$ trivial.
\end{corollary}

\begin{theorem}[Local structure]\label{thm:localstructure}
Every regular singular stratified bundle on $K$ satisfies
$\cE\cong\bigoplus_{c\in\Zp}\cL_c^{\oplus n_c}$ (finite sum).
Consequently $\Strat^{\rs}(K)$ is semisimple and equivalent to
$\Rep(\cT_p)$, $\cT_p:=\Spec k[\Zp/\ZZ]$; the $r$-spectrum modulo $\ZZ$ is
a complete invariant; and all $\Ext^1$ vanish.
\end{theorem}
\begin{proof}
If the $r$-spectrum is a single $c$ with multiplicity $N$, then
$\cL_{-c}\otimes\cE$ has $r$-spectrum $\{0\}$ by
\Cref{thm:functoriality}\eqref{it:additive}, hence is trivial by
\Cref{cor:allzero}.  In general the eigendecomposition of
\Cref{def:rspec} lifts to a filtration by $\cD(\log)$-stable saturated
subbundles (see \Cref{prop:blocks} for the lattice-level statement, whose
proof is independent of this theorem); the graded pieces are as above, and
$\Hom(\cL_c,\cL_{c'})=0$ for $c\not\equiv c'$ while
$\Ext^1(\cL_c,\cL_c)=\Ext^1(\cO,\cO)=0$ by \Cref{cor:allzero}.  Hence the
filtration splits.  Finally $\Zp/\ZZ$ has no $p$-torsion (its torsion is
$\ZZ_{(p)}/\ZZ$, of order prime to $p$), so $\cT_p$ is diagonalisable and
its representations are semisimple.
\end{proof}

\begin{corollary}\label{cor:nounipotent}
Local monodromy of a regular singular stratified bundle is semisimple.
There is no nilpotent element attached to $r$, hence no Jacobson--Morozov
$\mathfrak{sl}_2$ and no weight filtration built from local data.
\end{corollary}

\begin{example}[Non-semisimplicity is global]\label{ex:globalext}
On $\GG_m$ take $\sigma_n=\left(\begin{smallmatrix}1&t\\0&1\end{smallmatrix}\right)$.
Splitting needs $g_n-g_{n+1}^p=t$ in $k[t,t^{-1}]$; degrees satisfy
$d_n=\max(1,pd_{n+1})$, which is impossible.  So the extension is
non-split although $r=0$.  Over $k[[t]]$, $g=\sum_{i\ge0}t^{p^i}$ solves
$g=g^p+t$ and it splits, as \Cref{thm:localstructure} demands.
\end{example}

%% ===============================================================
\section{Comparison with Hecke invariants}\label{sec:comparison}
%% ===============================================================

Put $\tau_n:=\sigma_0\circ F^*\sigma_1\circ\cdots\circ F^{(n-1)*}\sigma_{n-1}$
and $\Sigma_n:=\im\tau_n$, so $\Sigma_0=\bar E_0$ and $\Sigma_n$ is spanned
by level-$n$ horizontal vectors; $\inv(\Sigma_n,\Sigma_{n+1})=p^n\mu_n$.

\begin{proposition}\label{prop:flagres}
If $\tau_n$ is upper triangular with diagonal $t^{A^{(1)}_n},\dots,t^{A^{(N)}_n}$
then for $m<n$, $\theta_{p^m}$ preserves the flag and
$\theta_{p^m}(e_j)\equiv\binom{-A^{(j)}_n}{p^m}e_j$ modulo lower terms and
$t\bar E_0$.
\end{proposition}
\begin{proof}
Let $f_i=\tau_n(\cdot)$, so $\partial^{[j]}f_i=0$ for $0<j<p^n$; then
$e_j=t^{-A_n^{(j)}}f_j+\sum_{i<j}c_{ij}f_i$ and the divided-power Leibniz
rule collapses to $\partial^{[p^m]}(cf)=\partial^{[p^m]}(c)f$.
\end{proof}

\begin{theorem}\label{thm:comparison}
The flag-wise (Iwasawa) invariants $\nu^{\Iw}_n:=(A^{(j)}_n)_j\in\Xst$
satisfy $\nu^{\Iw}_{n+1}=\nu^{\Iw}_n+p^n\mu^{\Iw}_n$, hence converge in
$\Xst\otimes\Zp$, and $-\lim_n\nu^{\Iw}_n=r$.
\end{theorem}

\begin{example}[Cartan invariants fail]\label{ex:cartanfails}
$G=\mathrm{GL}_2$, $g_0=\mathrm{diag}(t,t^{-1})$,
$g_1=\mathrm{diag}(t^{-1},t)$, $g_n=1$ for $n\ge2$.  Dominant
$\mu_0=\mu_1=(1,-1)$ so $\sum\mu_np^n=(1+p,-1-p)$, whereas
$g_0g_1^{(p)}=\mathrm{diag}(t^{1-p},t^{p-1})$ gives $\nu_n=(p-1,1-p)$ for
$n\ge2$ and $r=(p-1,1-p)$.  The discrepancy is the positive coroot
correction in $\inv(\Sigma_0,\Sigma_2)\le\mu_0+p\mu_1$, not divisible by
$p$.
\end{example}

\begin{example}[Convergence does not imply regularity]\label{ex:irregular}
$\sigma_n=\left(\begin{smallmatrix}1&t^{-1}\\0&1\end{smallmatrix}\right)$
gives $\nu_n=(-p^{n-1},p^{n-1})\to0$, yet splitting needs
$h_{n+1}^p-h_n=t^{-1}$, impossible since pole orders satisfy
$d_n=\max(pd_{n+1},1)$.  By \Cref{thm:localstructure} this bundle is not
regular singular: the logarithmic hypothesis cannot be dropped.
\end{example}

%% ===============================================================
\section{Rank one and low rank}\label{sec:examples}
%% ===============================================================

\begin{proposition}\label{prop:Gm}
$\Pic^{\strat}(\GG_m)\cong\Zp/\ZZ$.
\end{proposition}
\begin{proof}
Write $\sigma_n(1\otimes e_{n+1})=c_nt^{a_n}$; changes of basis act by
$a_n\mapsto a_n+pb_{n+1}-b_n$ and normalise $c_n=1$ as $k=\bar k$.  Set
$r=-\sum a_np^n$; then $r'-r=b_0-\lim b_Np^N=b_0$.  Surjectivity is the
digit expansion; injectivity: if $\sum a_np^n=0$ then $b_n:=A_n/p^n\in\ZZ$
works, $A_n=\sum_{i<n}a_ip^i$.
\end{proof}

\begin{proposition}[$p$-adic residue theorem]\label{prop:restheorem}
On $\PP^1$ with $D=\{0,\infty\}$: $r_0+r_\infty=-\deg\bar L_0$.
More generally $\sum_i\tr(r_i)[D_i]=-c_1(\bar E_0)$ in
$\mathrm{CH}^1(\bar X)\otimes\Zp$.
\CHECK{higher-dimensional case}
\end{proposition}

\begin{proposition}\label{prop:kummer}
$\bar r=0$ iff $\cE$ is trivial; $\bar r$ is torsion, i.e.\
$r\in\ZZ_{(p)}$, iff $\cE$ is a Kummer sheaf, using
$\ZZ_{(p)}/\ZZ\cong(\QQ/\ZZ)^{(p')}\cong\Hom(\pi_1^{\tame}(\GG_m),k^\times)$;
otherwise $\cE$ is a $p$-adic limit of Kummer sheaves not arising from
$\pi_1^{\tame}$.
\end{proposition}

\begin{corollary}\label{cor:nosection}
$\Zp\to\Zp/\ZZ$ has no continuous section, so there is no canonical
extension in general; on the torsion part $[0,1)\cap\ZZ_{(p)}$ is a
canonical section, which is exactly the range where Katz's canonical
extension exists.
\end{corollary}

\begin{example}[Irreducible non-tame rank two]\label{ex:GL2}
$p\ne2$, $f\colon\PP^1_z\to\PP^1_t$, $t=z^2$.  For $\cM$ of rank one with
residues $\alpha,\beta,\gamma,\delta$ at $0,1,-1,\infty$, the bundle
$\cE=f_*\cM$ on $\PP^1\setminus\{0,1,\infty\}$ has
\[
   r_0=\{\tfrac\alpha2,\tfrac{\alpha+1}2\},\quad r_1=\{\beta,\gamma\},\quad
   r_\infty=\{\tfrac\delta2,\tfrac{\delta+1}2\},
\]
irreducible when $\cM\not\cong\iota^*\cM$.  Taking
$\alpha\notin\ZZ_{(p)}$ gives an irreducible regular singular stratified
bundle that is not tame; hence
$\ker\bigl(\pi_1^{\strat,\rs}(U)\to\pi_1^{\tame}(U)\bigr)\ne1$.
The residue theorem checks: $\sum\tr r_x=\alpha+\beta+\gamma+\delta+1$ and
$\det f_*\cO=\cO(-1)$.
\end{example}

%% ===============================================================
\section{Boundary restriction}\label{sec:boundary}
%% ===============================================================

This section replaces \S8 of the first draft.  Fix a smooth irreducible
component $D_i$, let $N_i=\cO_{D_i}(D_i)$ and $D^{(i)}=\sum_{j\ne i}D_j$.
Write $\widehat\cO$ for the formal completion of $\cO_{\bar X}$ along
$D_i$; locally $\widehat\cO=\cO_{D_i}[[t]]$.

\subsection{Coordinate independence}
This is the point at which the first draft was incomplete.

\begin{proposition}\label{prop:coordfree}
The operators $R_a\in\End_{\cO_{D_i}}(\bar E_0|_{D_i})$, hence the
decomposition of \Cref{def:rspec}, are independent of the choice of local
equation $t$ for $D_i$.
\end{proposition}
\begin{proof}
The $D_a$ are the dual basis of the log divided-power coordinate on the
diagonal.  Precisely, let $p_1,p_2$ be the projections of the formal
logarithmic divided-power neighbourhood of the diagonal in
$\bar X\times\bar X$, put $t_j:=p_j^*t$ and
\[
   w:=\frac{t_2}{t_1}-1 ,
\]
which is a well-defined element of the log-PD-envelope with divided powers
$w^{[a]}$; the logarithmic stratification is
$\epsilon(1\otimes x)=\sum_aD_a(x)\otimes w^{[a]}$.

Let $t'=ut$ with $u$ a unit.  Then
$w'=\frac{t'_2}{t'_1}-1=\frac{u_2}{u_1}(1+w)-1$.  Now
$u_2-u_1$ lies in the ideal of the diagonal, and modulo the ideal
generated by $t_1$ together with the conormal directions along $D_i$ we
have $t_2-t_1=t_1w\equiv0$ and $u_2-u_1\equiv0$.  Restricting the
stratification to $D_i$ --- which is exactly the operation computing the
residues --- therefore gives $w'=w$, hence $D'_a=D_a$ on
$\bar E_0|_{D_i}$.
\end{proof}

\begin{remark}
Equivalently: the residue is the action of the \emph{log inertia group
scheme} $I_{D_i}$, the Cartier dual of the constant group $\Zp$, whose
representation category is that of $\cT_p$ from
\Cref{thm:localstructure}.  \Cref{prop:coordfree} says $I_{D_i}$ is
intrinsic.  Concretely one checks
$D'_a-D_a\in t\,\cD(\log D)+\cD(\log D)\cdot\mathrm{Der}_{D_i}$.
\end{remark}

\begin{corollary}\label{cor:globaldecomp}
$\bar E_0|_{D_i}=\bigoplus_{c\in\Zp}(\bar E_0|_{D_i})^c$ is a canonical
decomposition into $\cO_{D_i}$-subbundles, defined globally on $D_i$.
\end{corollary}
\begin{proof}
The $R_a$ are global $\cO_{D_i}$-linear endomorphisms
(\Cref{lem:leibniz,prop:coordfree}) with $R_a^p=R_a$, so the spectral
idempotents $\prod_{c'\ne c}(R_a-c')/(c-c')$ are global; multiplicities are
constant by \Cref{lem:constant}.
\end{proof}

\subsection{Residue jumps are Hecke exponents}

\begin{proposition}\label{prop:jump}
Let $M\subseteq L$ be $\cD(\log D_i)$-stable lattices in $E$ over
$\widehat\cO$ with $L/M$ supported on $D_i$.  If $L$ has $r$-spectrum
$\{c\}$ (a single value) and $M$ has $r$-spectrum $\{c'\}$, then
$a:=c'-c\in\ZZ_{\ge0}$ and $M=t^aL$.
\end{proposition}
\begin{proof}
For $x$ a $\binom{c}{\bullet}$-eigenvector, \Cref{lem:leibniz} gives
\[
   \theta_b(t^ax)=\sum_{i+j=b}D_i(t^a)\theta_j(x)
   =\sum_{i+j=b}\tbinom ai\tbinom cj\,t^ax=\tbinom{a+c}{b}t^ax
\]
by Vandermonde, so $t^aL$ has $r$-spectrum $\{c+a\}$; thus $c'-c$ is the
elementary divisor, in particular an integer.  If all elementary divisors
of $M\subseteq L$ equal $a$ then $t^aL\subseteq M$ and
$\mathrm{length}(L/M)=\mathrm{length}(L/t^aL)$, so $M=t^aL$.
\end{proof}

\begin{corollary}\label{cor:heckeexp}
Let $\{\bar E_n\}$ be a logarithmic lattice system and
$M_n:=\sigma_n(F^*\bar E_{n+1})$.  Then $M_n$ is $\cD(\log D_i)$-stable,
and the $r$-spectrum of $F^*\bar E_{n+1}$ is $p\cdot r(\bar E_{n+1})$.
Hence the Hecke exponents of $\sigma_n$ are
\[
   a_n=p\,r_{n+1}-r_n\ \in\ \ZZ^N ,\qquad r_n:=r(\bar E_n).
\]
\end{corollary}
\begin{proof}
$\cD(\log)$-stability: $\partial^{[a]}(1\otimes x)=0$ unless $p\mid a$ and
$D_{pb}(1\otimes x)=t^{pb}\otimes\partial^{[b]}x=1\otimes D_b(x)$.  Thus on
$F^*\bar E_{n+1}$ the residues are $R_a=0$ for $p\nmid a$ and
$R_{pb}=R_b(\bar E_{n+1})$: the digits are shifted up by one, i.e.\
$r(F^*\bar E_{n+1})=p\,r_{n+1}$.  Now apply \Cref{prop:jump}.
\end{proof}

\subsection{Block decomposition of the lattices}

\begin{proposition}\label{prop:blocks}
Let $\bar E$ be a logarithmic lattice over $\widehat\cO$.  Then
\[
   \bar E=\bigoplus_{\bar c\in\Zp/\ZZ}\bar E\{\bar c\}
\]
canonically, where each $\bar E\{\bar c\}$ is a $\cD(\log D_i)$-stable
$\widehat\cO$-subbundle whose restriction to $D_i$ is the sum of the
components $(\bar E|_{D_i})^c$ with $c\mapsto\bar c$.  The decomposition is
preserved by every morphism of logarithmic lattices, in particular by the
$\sigma_n$.
\end{proposition}
\begin{proof}
By \Cref{cor:rigidity} the commuting $k$-linear operators $\theta_a$ on
$\bar E$ are simultaneously diagonalisable with joint characters
$\binom{c}{\bullet}$, $c\in\Zp$; write $\bar E=\bigoplus_{c}\bar E[c]$ as
$k$-vector spaces.  Each $\theta_a$ is $\cO_{D_i}$-linear
(\Cref{lem:leibniz}), so each $\bar E[c]$ is an $\cO_{D_i}$-module, and by
the computation in \Cref{prop:jump},
\[
   t\cdot\bar E[c]\subseteq\bar E[c+1].
\]
Therefore $\bar E\{\bar c\}:=\bigoplus_{c\mapsto\bar c}\bar E[c]$ is
$t$-stable, i.e.\ an $\widehat\cO$-submodule, and the sum over
$\bar c\in\Zp/\ZZ$ is direct and exhausts $\bar E$.  Being a direct
summand of a projective module it is projective, hence a subbundle.
Morphisms of logarithmic lattices commute with the $\theta_a$ and so
preserve the $\bar E[c]$, a fortiori the blocks.  Finally, the blocks are
determined by the restriction to $D_i$ (\Cref{cor:globaldecomp}), which is
coordinate-free by \Cref{prop:coordfree}; hence so is the decomposition.
\end{proof}

\begin{remark}
The finer decomposition by $c\in\Zp$ does \emph{not} lift to the lattice:
residues differing by a positive integer are resonant, and the
corresponding extension of lattices need not split.  Grouping by
$\Zp/\ZZ$ is precisely what removes resonance.  Over $K$ resonance is
invisible ($\cL_c\cong\cL_{c'}$ iff $c-c'\in\ZZ$), which is why
\Cref{thm:localstructure} gives a full splitting there.
\end{remark}

\begin{lemma}[Normalisation]\label{lem:normalise}
Within a block $\bar E_n\{\bar c\}$ one may perform elementary
modifications along the subbundles $(\bar E_n|_{D_i})^{c'}$ so that the
$r$-spectrum of the block becomes a single value $c_n\in\Zp$ with
$c_n\mapsto\bar c$.  This may be done independently at each level.
\end{lemma}
\begin{proof}
An elementary modification $\ker(\bar E\to Q)$ along a subbundle $Q$ of
$\bar E|_{D_i}$ shifts the residues in $Q$ by $+1$ and leaves the others
fixed (\Cref{prop:jump}); the eigenvalues within a block differ by
integers, so finitely many such modifications equalise them.  The
$\sigma_n$ impose no constraint: for any choice of lattices,
$\sigma_n(F^*\bar E_{n+1})$ and $\bar E_n$ are commensurable lattices in
the same $K$-space.
\end{proof}

\subsection{The boundary restriction theorem}

\begin{theorem}[Boundary restriction]\label{thm:boundary}
Let $\cE$ be regular singular with logarithmic lattice system, normalised
as in \Cref{lem:normalise}, and fix a block $\bar c\in\Zp/\ZZ$ with
level-$n$ residue value $c_n\in\Zp$.  Put $B_n:=\bar E_n\{\bar c\}|_{D_i}$
and $a_n:=pc_{n+1}-c_n\in\ZZ$.  Then:
\begin{enumerate}
\item $\sigma_n$ induces isomorphisms
\[
   F_{D_i}^*\,B_{n+1}\ \isoto\ B_n\otimes N_i^{\otimes(-a_n)} ,
\]
and each $B_n$ carries a logarithmic connection along
$(D_i,D_i\cap D^{(i)})$ compatible with these.
\item Setting $e_n:=c_n$ solves the untwisting recursion
      $e_n=pe_{n+1}-a_n$; thus $\bigl(B_n\bigr)$ is an
      \emph{$N_i$-twisted regular singular stratified bundle of weight
      $c:=c_0$} on $(D_i,D_i\cap D^{(i)})$, with total twist
      ``$N_i^{\otimes c}$''.
\item $\bigl(B_n\bigr)$ becomes an honest regular singular stratified
      bundle after tensoring by a line bundle if and only if there exist
      $M_n\in\Pic(D_i)$ with $[M_n]=p[M_{n+1}]-a_n[N_i]$, i.e.\ iff the
      element $c\cdot[N_i]$ of the $p$-adic completion
      $\widehat{\Pic(D_i)}_p$ lies in the image of $\Pic(D_i)$.  The
      obstruction thus lies in $\widehat{\Pic(D_i)}_p/\Pic(D_i)$, and
      vanishes when $c\in\ZZ$ or $N_i$ is torsion.
\item The adjoint tractor bundle $\cA|_{D_i}$ has cancelling twists and is
      always an honest log parabolic geometry on
      $(D_i,D_i\cap D^{(i)})$.
\item $\cR_i\colon\cE\mapsto\bigl(\bar E_\bullet\{\bar c\}|_{D_i}\bigr)_{\bar c}$
      is compatible with $\otimes$, duals and pullbacks, and residues along
      $D^{(i)}$ restrict as in
      \Cref{thm:functoriality}\eqref{it:pullback}.
\end{enumerate}
\end{theorem}
\begin{proof}
(1) By \Cref{prop:blocks} the $\sigma_n$ preserve blocks; after
normalisation each block of $\bar E_n$ has a single residue value, so
\Cref{cor:heckeexp} and \Cref{prop:jump} give
$\sigma_n(F^*\bar E_{n+1}\{\bar c\})=t^{a_n}\bar E_n\{\bar c\}
=\bar E_n\{\bar c\}(-a_nD_i)$.  Restricting to $D_i$ and using
$\bigl(F^*\bar E_{n+1}\bigr)\big|_{D_i}=F_{D_i}^*\bigl(\bar E_{n+1}|_{D_i}\bigr)$
(absolute Frobenius preserves $D_i$) and
$\cO(-a_nD_i)|_{D_i}=N_i^{\otimes(-a_n)}$ gives the displayed isomorphism.
Horizontality along $D_i$: a coordinate $s$ along $D_i$ satisfies
$[\,\delta^{(m)}_t,\partial_s^{[j]}\,]=0$, so all the constructions above
are flat in the $D_i$-directions, and the logarithmic poles along
$D_j$, $j\ne i$, are preserved.

(2) $pc_{n+1}-a_n=pc_{n+1}-pc_{n+1}+c_n=c_n$.

(3) Untwisting means finding $M_n$ with $F^*(B_{n+1}\otimes M_{n+1})\cong
B_n\otimes M_n$, i.e.\ $[M_n]=p[M_{n+1}]-a_n[N_i]$.  Iterating,
$[M_0]=-\bigl(\sum_{i<n}a_ip^i\bigr)[N_i]+p^n[M_n]$, so $[M_0]$ must
reduce to $-\bigl(\sum_{i<n}a_ip^i\bigr)[N_i]$ modulo $p^n\Pic(D_i)$ for
all $n$; since $\sum_na_np^n=-c$ this says $[M_0]\mapsto c\cdot[N_i]$ in
$\widehat{\Pic(D_i)}_p$.  Conversely such an $M_0$ determines the $M_n$.

(4) and (5) are formal.
\end{proof}

\begin{corollary}[Extension criterion]\label{cor:extension}
$\cE$ extends to a stratified bundle in a neighbourhood of $D_i$ iff
$r_i\in\Xst$.
\end{corollary}
\begin{proof}
$(\Leftarrow)$ Normalise all residues to $0$ by \Cref{lem:normalise}; then
$a_n=0$, so every $\sigma_n$ is a lattice isomorphism and
$(\bar E_n,\sigma_n)$ is an $F$-divided sheaf near $D_i$.  $(\Rightarrow)$
clear.
\end{proof}

\begin{example}
$\bar X=\PP^n$, $D_i$ a smooth hypersurface of degree $d$, so
$\Pic(D_i)=\ZZ$ (for $n\ge4$) and $[N_i]=d$.  \Cref{thm:boundary}(3) then
says the boundary object untwists iff $cd\in\ZZ$, i.e.\ iff
$c\in\frac1d\ZZ$ --- a strong constraint recovering
\Cref{cor:extension} when $d=1$.
\end{example}

\subsection{The dictionary with conformal infinity}\label{sec:conformal}

\begin{center}
\begin{tabular}{@{}p{0.44\textwidth}p{0.48\textwidth}@{}}
\hline
\textbf{Characteristic $0$ (Poincar\'e--Einstein)} &
\textbf{Characteristic $p$}\\
\hline
conformal infinity $\partial M$ & boundary component $D_i$\\
defining function $\rho$ & logarithmic lattice $\bar E_0$\\
conformal class & indeterminacy of $r$ by $\Xst$ (\Cref{thm:welldef})\\
asymptotic exponents & $p$-adic digits of $r$\\
conformal density of weight $w$ & $N_i^{\otimes c}$-twist (\Cref{thm:boundary})\\
tractor bundle has weight $0$ & $\cA|_{D_i}$ untwisted\\
quasi-unipotent monodromy & torsion $r\in\ZZ_{(p)}$, i.e.\ tameness\\
unipotent part, weight filtration & absent locally
   (\Cref{cor:nounipotent}); global $\Ext$ only\\
\hline
\end{tabular}
\end{center}

The identification ``residue $=$ conformal weight'' explains the
$p$-adicity structurally: $N_i^{\otimes c}$ is meaningful in the
$F$-divided formalism exactly for $c\in\Zp$, because the recursion
$e_n=pe_{n+1}-a_n$ is a $p$-adic expansion.

%% ===============================================================
\section{$p$-adic middle convolution}\label{sec:MC}
%% ===============================================================

We now construct a middle convolution with parameter in $\Zp$.  The
essential input is \Cref{thm:localstructure}: because local monodromy is
semisimple, the local invariants entering the theory reduce to
multiplicities of residues, and all of Katz's Jordan-block bookkeeping
disappears.

\subsection{Convolution via Gauss--Manin}
Let $S=\{x_1,\dots,x_s\}\subset\AAA^1$ and $U=\AAA^1\setminus S$.  For
$c\in\Zp$ let $\cL_c$ be the rank one stratified bundle on $\GG_m$ with
residues $c$ at $0$ and $-c$ at $\infty$ (\Cref{prop:Gm}), and write
$\cL_c(y-x)$ for its pullback under $(x,y)\mapsto y-x$.

\begin{construction}\label{con:MC}
For $\cE\in\Strat^{\rs}(U)$ and $c\in\Zp$ put, for $y$ varying in
$\AAA^1\setminus S$,
\[
   \MC_c(\cE)_y:=\im\Bigl(
     H^1_{\dR,c}\bigl(\AAA^1\setminus(S\cup\{y\}),\ \cE\otimes\cL_c(y-x)\bigr)
     \longrightarrow
     H^1_{\dR}\bigl(\cdots\bigr)\Bigr).
\]
Relative de Rham cohomology of a stratified bundle along a smooth
morphism carries a Gauss--Manin stratification (Katz--Oda; the
construction is compatible with $F$-divided systems), so $\MC_c(\cE)$ is
a stratified bundle on $\AAA^1\setminus S$ wherever the cohomology sheaves
are locally free of the expected rank.
\end{construction}

\begin{remark}[Why the middle extension is unambiguous]\label{rem:middle}
In characteristic $0$ the definition of $j_{!*}$ requires bookkeeping of
the unipotent part of local monodromy: the local invariants
$\dim\cF^{I_x}$ is the \emph{number of Jordan blocks} with eigenvalue $1$,
not the multiplicity of the eigenvalue.  By \Cref{thm:localstructure} our
local monodromy is semisimple, so
\[
   \dim\cF^{I_x}=\#\{\,j:\ r_{x,j}\in\ZZ\,\}=:\delta_x ,
\]
the plain multiplicity.  Note that \Cref{cor:nosection} forbids a
canonical extension in the non-torsion range, but this is irrelevant here:
the cohomological definition of $\MC_c$ requires no normalisation.  This
is the concrete payoff of the local structure theorem.
\end{remark}

\subsection{Rank and exponents}

\begin{proposition}\label{prop:MCrank}
Assume $c\notin\ZZ$ and $\cE$ irreducible and non-trivial with
$r$-spectra $r_i$ at $x_i$ and $r_\infty$ at $\infty$, $N=\rk\cE$.  Then
\[
   \rk\MC_c(\cE)=Ns-\sum_{i=1}^s\delta_i-\delta_\infty(c),
   \qquad
   \delta_i:=\#\{j: r_{i,j}\in\ZZ\},\quad
   \delta_\infty(c):=\#\{j: r_{\infty,j}-c\in\ZZ\}.
\]
\end{proposition}
\begin{proof}
For a regular singular $\cF$ of rank $M$ on $\PP^1$ with singular set $T$
and $j\colon\PP^1\setminus T\hookrightarrow\PP^1$, one has
$\chi(\PP^1,j_!\cF)=M(2-|T|)$ and
$\chi(\PP^1,j_{!*}\cF)=M(2-|T|)+\sum_{x\in T}\dim\cF^{I_x}$, so
$\dim H^1_{\mathrm{mid}}=M(|T|-2)-\sum_x\dim\cF^{I_x}$.  Apply this to
$\cF=\cE\otimes\cL_c(y-x)$ with $T=\{x_1,\dots,x_s,y,\infty\}$,
$|T|=s+2$; the residues are $r_i$ at $x_i$, $c$ at $y$ (so
$\dim\cF^{I_y}=0$ as $c\notin\ZZ$) and $r_{\infty,j}-c$ at $\infty$.  Use
\Cref{rem:middle} to identify $\dim\cF^{I_x}$ with $\delta_x$.
\end{proof}

\begin{proposition}[Exponent rules]\label{prop:MCrules}
In the situation of \Cref{prop:MCrank}, write $N'=\rk\MC_c(\cE)$.  The
$\Zp$-valued Riemann scheme of $\MC_c(\cE)$ is
\[
\begin{aligned}
   \text{at }x_i:&\quad \{\,r_{i,j}+c\ :\ r_{i,j}\notin\ZZ\,\}
     \ \cup\ \{\,0\ \text{with multiplicity}\ N'-(N-\delta_i)\,\},\\[2pt]
   \text{at }\infty:&\quad \{\,r_{\infty,j}-c\ :\ r_{\infty,j}-c\notin\ZZ\,\}
     \ \cup\ \{\,-c\ \text{with multiplicity}\ N'-(N-\delta_\infty(c))\,\},
\end{aligned}
\]
all equalities modulo $\ZZ$, i.e.\ up to the choice of lattice.
\end{proposition}

\begin{remark}
These are Katz's rules \cite[Ch.~6]{Katz2} transported to $\Zp$ by
$\mu=e^{2\pi i\alpha}\rightsquigarrow\alpha$, with the simplification of
\Cref{rem:middle}: the rules are now purely additive in $\Zp$, with no
Jordan-block clause.
\CHECK{a direct proof via the local behaviour of the Gauss--Manin
connection, rather than transport from characteristic $0$, should be
written out}
\end{remark}

\begin{example}[$p$-adic Gauss hypergeometric]\label{ex:2F1}
Let $\cE=\cL_{\alpha}(x)\otimes\cL_{\beta}(x-1)$ on
$\PP^1\setminus\{0,1,\infty\}$: rank $1$, residues $\alpha$ at $0$,
$\beta$ at $1$, $-\alpha-\beta$ at $\infty$; assume
$\alpha,\beta,\alpha+\beta+c\notin\ZZ$.  Then $s=2$,
$\delta_1=\delta_2=\delta_\infty(c)=0$, so
\[
   \rk\MC_c(\cE)=1\cdot2-0-0=2,
\]
with Riemann scheme
\[
   \begin{array}{c|c|c}
   0 & 1 & \infty\\\hline
   0,\ \alpha+c & 0,\ \beta+c & -c,\ -\alpha-\beta-c
   \end{array}
\]
Comparing with the classical scheme
$\bigl(\begin{smallmatrix}0,1-\gamma&0,\gamma-a-b&a,b\end{smallmatrix}\bigr)$
we get $\gamma=1-\alpha-c$, $a+b=1-\alpha-\beta-2c$, $ab$ determined; the
sum of all six exponents is $0$, consistent with
\Cref{prop:restheorem} and $\deg\bar E_0=0$ (the classical value $1$ is a
different lattice normalisation).  We call the resulting object the
\emph{$p$-adic hypergeometric stratified bundle}
$\cH(a,b;\gamma)$, $a,b,\gamma\in\Zp$.
\end{example}

\begin{remark}[Relation to Dwork]
A $\Zp$-parameter hypergeometric is, level by level, a compatible system
of hypergeometric-type equations with parameters
$a,b,\gamma\bmod p^{n+1}$.  This is precisely the shape of Dwork's
$p$-adic hypergeometric theory and of his congruences for
${}_2F_1$-coefficients; \Cref{ex:2F1} should be compared with the
Frobenius structure there.
\TODO{make the comparison with Dwork precise}
\end{remark}

\begin{example}[Consistency with \Cref{ex:GL2}]
Taking $c=\tfrac12$ (possible for $p\ne2$) and $\alpha$ suitable, the
scheme of \Cref{ex:2F1} reproduces the exponent pattern
$\{\tfrac\alpha2,\tfrac{\alpha+1}2\}$ of the pushforward along the double
cover, as it must.
\end{example}

\subsection{Rigidity and the $p$-adic Katz algorithm}

\begin{definition}
For $\cE$ regular singular on $\PP^1\setminus T$ define the
\emph{rigidity index}
\[
   \rig(\cE):=2N^2-\sum_{x\in T}\bigl(N^2-\dim Z(r_x)\bigr),
   \qquad
   \dim Z(r_x)=\sum_{\bar c\in\Zp/\ZZ}m_{x,\bar c}^2 ,
\]
where $m_{x,\bar c}$ is the multiplicity of $\bar c$ in the $r$-spectrum
at $x$.  By \Cref{thm:localstructure} the centraliser of local monodromy
really has this dimension --- no Jordan blocks intervene.  $\cE$ is
\emph{rigid} if $\rig(\cE)=2$.
\end{definition}

\begin{proposition}\label{prop:rigMC}
$\rig(\MC_c(\cE))=\rig(\cE)$ for $c\notin\ZZ$, and $\MC_c$ preserves
irreducibility.  Moreover $\MC_{c'}\circ\MC_c=\MC_{c+c'}$ on irreducible
non-trivial objects, and $\MC_{-c}\circ\MC_c=\id$.
\CHECK{transport of \cite[6.0]{Katz2}; needs the direct Gauss--Manin proof}
\end{proposition}

\begin{conjecture}[$p$-adic Katz algorithm]\label{conj:katzalg}
Every irreducible rigid regular singular stratified bundle on
$\PP^1\setminus T$ is obtained from a rank one object by a finite sequence
of operations $\cE\mapsto\cE\otimes\cL$ ($\cL$ of rank one) and
$\cE\mapsto\MC_c(\cE)$, $c\in\Zp$.
\end{conjecture}

\begin{remark}
The choice of $c$ in Katz's reduction step is $c=-\sum(\text{certain
residues})$.  In characteristic $0$ this requires a root of unity of the
right order in the coefficient field; here $c$ ranges over the ring $\Zp$
and such sums always exist, so the algorithm has \emph{more} room, not
less.  The condition to check at each step is only $c\notin\ZZ$.
\end{remark}

\subsection{What is still missing}

\begin{problem}\label{prob:MCrs}
Does $\MC_c$ preserve regular singularity?  Equivalently: is the
Gauss--Manin stratification of a regular singular stratified bundle along
a family of affine curves again regular singular?  In characteristic $0$
this is standard; here it is open, and it is the main gap in
\Cref{sec:MC}.
\end{problem}

\begin{problem}
Local freeness and base change for $H^1_{\dR}$ in
\Cref{con:MC}: identify hypotheses (e.g.\ on $\delta_x$) guaranteeing that
$\MC_c(\cE)$ is a vector bundle of the rank predicted by
\Cref{prop:MCrank}.
\end{problem}

\begin{problem}
Give an intrinsic proof of \Cref{prop:MCrules} by computing the local
Gauss--Manin residue, and deduce \Cref{prop:rigMC}.
\end{problem}

\begin{problem}
Combine $\MC_c$ with \Cref{thm:boundary}: is there a middle convolution in
higher dimension compatible with boundary restriction, so that the
$p$-adic Katz algorithm can be run inductively on $\dim\bar X$?
\end{problem}

%% ===============================================================
\section{Consequences and the inductive framework}\label{sec:applications}
%% ===============================================================

\begin{corollary}\label{cor:blocks2}
$\Strat^{\rs}(U)$ decomposes into blocks indexed by
$\prod_i(\Xst\otimes\Zp/\Xst)/W$; tame and non-tame blocks are separated.
\end{corollary}

\begin{proposition}
Since local monodromy is always semisimple, all non-semisimplicity of
$\pi_1^{\strat,\rs}(U)$ is global; one expects
\[
   1\to(\text{interior})\to\pi_1^{\strat,\rs}(U)\to\prod_i\cT_p,
\]
with the tame quotient factoring through $\ZZ_{(p)}/\ZZ$.
\TODO{the log Gieseker--Katz statement}
\end{proposition}

\subsection{Programme}
\Cref{thm:boundary} is dimension-decreasing, so induction on
$\dim\bar X$ is available, provided the inductive hypothesis is stated for
$N_i$-twisted objects, carrying $\widehat{\Pic(D_i)}_p$ along.  Targets:
\begin{enumerate}
\item $\pi_1^{\strat,\rs}(\AAA^n)=1$ and Lefschetz-type theorems, using
      \Cref{cor:extension} and
      \Cref{thm:functoriality}\eqref{it:pullback} with $e=1$;
\item the residue theorem in $\mathrm{CH}^1\otimes\Zp$ and non-existence
      results;
\item a new proof of the tame Gieseker conjecture by induction on
      dimension;
\item the $p$-adic Katz algorithm (\Cref{conj:katzalg}) and a
      classification of rigid objects on $\PP^1$;
\item twisted induction: the correct formulation of the inductive
      hypothesis for $N_i^{\otimes c}$-twisted stratified bundles.
\end{enumerate}

%% ===============================================================
\section{Open problems}
%% ===============================================================

\begin{problem}
Is $\Ext^1$ of two objects with non-congruent residues always zero over
$U$ as well, or only locally?  Describe the global $\Ext$-object
$N\in\varprojlim_n\Ext^1(\gr,\gr)$ as the replacement for the unipotent
part.
\end{problem}

\begin{problem}
Compare $\cT_p=\Spec k[\Zp/\ZZ]$ with the full local stratified
fundamental group of $K$, including irregular objects
(\Cref{ex:irregular}).
\end{problem}

\begin{problem}
Resolve \Cref{prob:MCrs}, the preservation of regular singularity by
$\MC_c$.
\end{problem}

%% ===============================================================
\begin{thebibliography}{99}

\bibitem{Dwork} B.~Dwork,
  \emph{$p$-adic cycles}, Publ. Math. IHES \textbf{37} (1969), 27--115.
  \TODO{also cite Dwork's hypergeometric congruences}

\bibitem{EsnaultKindler} H.~Esnault and L.~Kindler,
  \emph{Lefschetz theorems for tamely ramified coverings}.
  \TODO{complete reference}

\bibitem{Gieseker} D.~Gieseker,
  \emph{Flat vector bundles and the fundamental group in non-zero
  characteristics}, Ann. Scuola Norm. Sup. Pisa \textbf{2} (1975), 1--31.

\bibitem{Katz} N.~Katz,
  \emph{Nilpotent connections and the monodromy theorem: applications of a
  result of Turrittin}, Publ. Math. IHES \textbf{39} (1970), 175--232.

\bibitem{Katz2} N.~Katz,
  \emph{Rigid Local Systems}, Ann. of Math. Studies \textbf{139},
  Princeton, 1996.

\bibitem{Kindler} L.~Kindler,
  \emph{Regular singular stratified bundles and tame ramification}.
  \TODO{complete reference}

\bibitem{DettweilerReiter} M.~Dettweiler and S.~Reiter,
  \emph{An algorithm of Katz and its application to the inverse Galois
  problem}, J. Symbolic Comput. \textbf{30} (2000), 761--798.

\bibitem{dosSantos} J.~P.~P. dos Santos,
  \emph{Fundamental group schemes for stratified sheaves},
  J. Algebra \textbf{317} (2007), 691--713.

\end{thebibliography}

\end{document}
